\section{Background}
\label{sec:app:background}

% \begin{definition}[Sub-differential]
%     For a convex function $f: \mathcal{X} \to \mathbb{R}^*$, the sub-differential $\partial f(x)$ is defined as 
%     $$
%     \partial f(x) = \{ y \in \mathcal{X}: \forall z, f(x + z) \geq f(x) + \langle y, z \rangle \}
%     $$
% \end{definition}

% \begin{definition}[Fenchel conjugate]
%     The Fenchel conjugate, $f^*: \mathcal{X} \to \mathbb{R}^*$ for a convex function $f: \mathcal{X} \to \mathbb{R}^*$ is defined as $f^*(y) = \sup\limits_{x \in \mathcal{X}} \langle x, y \rangle - f(x)$
% \end{definition}

% \begin{definition}[Dual norm]
%     Given a norm $\lVert \cdot \rVert$ on $\mathcal{X}$, the dual norm $\lVert \cdot \rVert_*$ is defined as 
%     $$
%     \lVert y \rVert_* = \sup \{ \langle x, y \rangle: \lVert x \rVert \leq 1 \}
%     $$
% \end{definition}

% Note that the Fenchel conjugate of $\frac{1}{2} \lVert x \rVert$ is $\frac{1}{2} \lVert y \rVert_*$.

\begin{definition}[Strong convexity]
    A function $f: \mathcal{X} \to \mathbb{R}^*$ is $\mu-$strongly convex with respect to $\lVert \cdot \rVert$ if $\forall x, y$ in the relative interior of $\text{dom } f$ and $\alpha \in (0, 1)$,
    $$
    f(\alpha x + (1 - \alpha) y) \leq \alpha f(x) + (1 - \alpha) f(y) - \frac{1}{2} \mu \alpha (1 - \alpha) \lVert x - y \rVert^2
    $$
\end{definition}

\begin{definition}[Smoothness]
    A function $f: \mathcal{X} \to \mathbb{R}$ is $\beta-$smooth with respect to $\lVert \cdot \rVert$ if $f \in C^1$ and if $\forall x, y \in \text{dom } f$,
    $$
    f(x + y) \leq f(x) + \langle \nabla f(x), y \rangle + \frac{1}{2} \beta \lVert y \rVert^2
    $$
\end{definition}

\subsection{The standard covering number guarantee}
\label{sec:covering}

We reiterate the standard bound under the covering argument from \cite{duchi2023} here. This bounds the deviation of the mini-batch approximation of the supremum from the true value.

\input{bound}

% \begin{definition}[Dense set]
%     A subset $A$ of a topological space $X$ is said to be \textit{dense} in $X$ if every item in $X$ is present in or is arbitrarily close to some point in $A$.
% \end{definition}

% \begin{definition}(Lebesgue integral)
%     The Lebesgue integral is a generalization of the Riemann integral. It is defined in terms of the Lebesgue sum
%     \[
%         \sum_i \alpha_i \mu(S_i)
%     \]
%     where $\alpha_i$ is the value of the function in the subinterval $i$, and $\mu(S_i)$ is the Lebesgue measure of the interval $S_i$ of points for which the value of the function is approximately $\alpha_i$. A function is \textit{Lebesgue integrable} if its Lebesgue integral is finite.
% \end{definition}

% \begin{definition}[$L^p$ space]
%     An $L^p$ space is a space of measurable functions for which the $p$th power is Lebesgue integrable.
% \end{definition}